\chapter{اللقاء السادس والعشرون: ختام أحكام الحج وبداية أوصاف المنافقين والمجاهدين}

\section{مقدمة}
السلام عليكم ورحمة الله وبركاته. الحمد لله رب العالمين، وصلى الله على محمد وعلى آله وصحبه وسلم تسليماً كثيراً.

استأنف الشيخ هذا اللقاء من قوله تعالى: \begin{ayah}لَيْسَ عَلَيْكُمْ جُنَاحٌ أَنْ تَبْتَغُوا فَضْلًا مِنْ رَبِّكُمْ\end{ayah}، ثم مضى في آيات الإفاضة من عرفات والمشعر الحرام وأيام الذكر والنفر، حتى انتهى إلى أول قوله تعالى: \begin{ayah}وَمِنَ النَّاسِ مَنْ يُعْجِبُكَ قَوْلُهُ فِي الْحَيَاةِ الدُّنْيَا\end{ayah} وما بعدها من المقابلة بين المنافق المفسد وبين من يبيع نفسه لله طلباً لمرضاته.

\separator

\section{تأويل (ليس عليكم جناح أن تبتغوا فضلاً من ربكم)}
\isnad{قال أبو جعفر:} لا حرج عليكم في التماس الرزق والتجارة في مواسم الحج قبل الإحرام وبعده، ولا ينافي ذلك صحة النسك ولا يقدح في الإخلاص إذا استقام القصد.

وذكر \rawi{الطبري} أن الآية نزلت في قوم كانوا يكرهون البيع والشراء في الحج، أو يظنون أن الاشتغال بالمكاسب يسقط عنهم اسم الحج، فرفع الله عنهم هذا الوهم، وأباح لهم أن يبتغوا فضلاً من ربهم مع إقامة المناسك.

\begin{fawaid}[المباحات لا تخرج العبادة عن حقيقتها إذا استقامت النية]
ليس كل انتفاع دنيوي يصاحب العبادة يقدح فيها. فمن أدى النسك على وجهه، ثم باع واشترى أو اكتسب رزقاً مباحاً، فلا يخرجه ذلك عن كونه حاجاً، ما دام قد أتى بما أمر الله به.
\end{fawaid}

\separator

\section{تأويل (فإذا أفضتم من عرفات)}
انتقل \rawi{الطبري} هنا إلى تفسير الإفاضة من عرفات، فتكلم أولاً على معنى الإفاضة في اللسان، ثم على اسم \textit{عرفات} من جهة بنائه العربي وسبب تسميته، وذكر في ذلك أقوال أهل اللغة وأهل العلم.

ونبّه الشيخ إلى دقة صنيع \rawi{الطبري} في التفريق بين أنواع الخلاف:
\begin{itemize}
    \item خلاف لغوي في بناء الكلمة وتصريفها.
    \item وخلاف علمي في سبب تسميتها ومعناها.
    \item وخلاف فقهي في حدود الموضع الذي يدخل في عرفات وما لا يدخل فيه.
\end{itemize}

ومن خلال هذا الموضع يظهر أن \rawi{الطبري} يرجع كل جزء من الآية إلى أهله: فإن كان الكلام في اللسان نقل عن النحاة واللغويين، وإن كان في التفسير والأثر رجع إلى أهل التأويل.

\begin{fawaid}[إرجاع كل مسألة إلى أهلها من تمام التحقيق]
من قوة تفسير الطبري أنه لا يخلط بين أبواب العلم؛ فإذا كانت المسألة لغوية رجع فيها إلى أهل العربية، وإذا كانت تفسيرية أو أثرية رجع فيها إلى أهل التأويل، ثم يزن الجميع بميزان الدليل والسياق.
\end{fawaid}

\section{تأويل (فاذكروا الله عند المشعر الحرام)}
فسر \rawi{الطبري} الآية بالأمر بذكر الله عند المشعر الحرام بعد الإفاضة من عرفات، وذكر ما جاء عن السلف في حد المشعر الحرام وما يستحب فيه من الذكر والدعاء والصلاة.

\section{تأويل (واذكروه كما هداكم وإن كنتم من قبله لمن الضالين)}
في هذا الموضع قرر \rawi{الطبري} أن الذكر المشروع هنا ليس مجرد ألفاظ تؤدى، بل هو شكر لله على نعمة الهداية، إذ نقلهم من حال الجاهلية والضلال إلى منسك مشروع مضبوط بالوحي.

\separator

\section{تأويل (ثم أفيضوا من حيث أفاض الناس)}
ذكر \rawi{الطبري} في هذه الآية قولين مشهورين:
\begin{itemize}
    \item أن الخطاب لقريش ومن كان على طريقتهم من الحمس، إذ كانوا يقفون دون عرفات تمييزاً لأنفسهم، فأمروا أن يفيضوا من حيث أفاض سائر الناس.
    \item وأن الخطاب عام للمسلمين، ومعنى الناس فيه إبراهيم ومن تبعه في سنته.
\end{itemize}

وظهر من مناقشة \rawi{الطبري} أنه يعتني هنا غاية العناية بالسياق؛ لأن بعض الأوجه اللغوية وإن أمكن تخريجها، إلا أن نظم الآية وسياقها يدفعها أو يضعفها.

\begin{fawaid}[السياق من أقوى ما يرد به على التخريجات اللغوية البعيدة]
كثيراً ما ينقل الطبري عن أهل اللغة وجوهاً محتملة، ثم يرد بعضها لا لضعفها من جهة اللسان فقط، بل لأن سياق الآية يأباها. فسلامة العربية وحدها لا تكفي إذا خالفت نظم الكلام واتجاهه.
\end{fawaid}

\section{تأويل (فإذا قضيتم مناسككم فاذكروا الله كذكركم آباءكم أو أشد ذكراً)}
فسر \rawi{الطبري} هذا الموضع بأن العرب كانوا إذا فرغوا من حجهم جلسوا يعددون مآثر آبائهم ومفاخرهم، فأمرهم الله أن يكون ذكرهم له بعد قضاء المناسك أعظم من ذكرهم لآبائهم، بل أشد وأخلص.

\section{تأويل (فمن الناس من يقول ربنا آتنا في الدنيا)}
بيّن \rawi{الطبري} أن هذا وصف لمن قصر همته على الدنيا، فلم يسأل الله إلا العاجلة، فحرم نفسه خير الآخرة.

\section{تأويل (ومنهم من يقول ربنا آتنا في الدنيا حسنة وفي الآخرة حسنة)}
في مقابل الصنف الأول جاء وصف المؤمن الجامع بين خير الدنيا والآخرة، وهو أعدل الطلب وأكمله؛ لأن الشريعة لا تقطع العبد عن مصالح دنياه، لكنها ترده إلى طلبها على وجه يعينه على الآخرة.

\separator

\section{تأويل (واذكروا الله في أيام معدودات فمن تعجل في يومين فلا إثم عليه ومن تأخر فلا إثم عليه)}
وقف \rawi{الطبري} في هذا الموضع وقفة مطولة، فذكر الأقوال في معنى نفي الإثم عن المتعجل والمتأخر، ثم ناقشها واحداً واحداً، وبيّن ما فيها من الضعف، ثم رجح ما يوافق ظاهر التنزيل والسنة الثابتة.

ونبّه الشيخ إلى أن هذا الموضع من أنفس المواضع لاكتشاف نسق \rawi{الطبري} الكامل في عرض الخلاف:
\begin{itemize}
    \item يذكر الأقوال.
    \item وينسبها إلى قائليها.
    \item ويبين حجة كل قول.
    \item ثم يخطئ ما يراه خطأً بدليل مفصل.
    \item وربما استعمل في ذلك طريقة الصبر والتقسيم.
\end{itemize}

\begin{fawaid}[من أكمل صور منهج الطبري أن يذكر القول ثم ينقضه بحجته]
لا يكتفي الطبري برد القول المخالف حكماً مجرداً، بل يبين مأخذه، ثم يكشف موضع خطئه، ويقيم الحجة على ذلك من ظاهر الآية أو السنة أو الإجماع أو حسن التقسيم العقلي. وهذا من أنفع ما يتربى عليه طالب العلم في باب النظر.
\end{fawaid}

\separator

\section{تأويل (واتقوا الله واعلموا أنكم إليه تحشرون)}
ختم \rawi{الطبري} هذا المقطع بالتذكير بأن أحكام الحج وآدابه وحدوده إنما تؤدى تحت رقابة التقوى، وأن المرجع إلى الله الذي يجازي المحسن بإحسانه والمسيء بإساءته.

\separator

\section{تأويل (ومن الناس من يعجبك قوله في الحياة الدنيا)}
\isnad{قال أبو جعفر:} هذا وصف لصنف من الناس حسن العبارة، معجب الظاهر، سيئ الباطن، يستشهد الله كذباً على ما في قلبه، وهو في حقيقة أمره شديد الخصومة، ممارٍ بالباطل، مخاتل في لسانه وقلبه.

ونقل \rawi{الطبري} في سبب النزول أقوالاً:
\begin{itemize}
    \item فقيل: نزلت في الأخنس بن شريق.
    \item وقيل: في طائفة من المنافقين تكلموا في بعض سرايا المسلمين.
    \item وقيل: هي عامة في كل من كانت سريرته تخالف علانيته.
\end{itemize}

ثم رجح المعنى العام مع عدم إهدار سبب النزول، ونبّه الشيخ إلى أن الصواب ليس في إهمال من نزلت فيه الآية، ولا في حبسها عليه وحده، بل في البدء به ثم توسيع الدلالة إلى كل من شمله ظاهرها.

\begin{fawaid}[البداية بسبب النزول لا تعني حبس الآية عليه]
من الخطأ أن نهمل الشخص أو الحادثة التي نزلت فيها الآية، كما أن من الخطأ أن نقصر حكمها عليه وحده. والمنهج المحكم أن نبدأ بسبب النزول لأنه يوضح السياق، ثم نعمم الدلالة بحسب ظاهر النص.
\end{fawaid}

\section{تأويل (ويشهد الله على ما في قلبه وهو ألد الخصام)}
تكلم \rawi{الطبري} على القراءتين في قوله \textit{ويشهد الله} و\textit{ويشهدُ الله}، ثم رجح قراءة الجمهور، وفسر \textit{ألد الخصام} بما يجمع شدة الجدل، والاعوجاج عن الحق، والخصومة بالباطل.

ونبّه الشيخ هنا إلى فائدة دقيقة من طريقة \rawi{الطبري}: فهو يصنف الأقوال، لكنه أحياناً يرجع بعضها إلى بعض إذا كانت متقاربة في الحقيقة، فلا يضخم الخلاف بلا موجب.

\begin{fawaid}[ليس كل اختلاف في العبارة اختلافاً حقيقياً في المعنى]
من دقة الطبري أنه قد يذكر قولين أو أكثر، ثم يبين أن بعضها متقارب المعنى، وأن الاعوجاج في الخصومة داخل في شدة الجدال، فلا يجعل كل تنوع لفظي خلافاً مستقلاً.
\end{fawaid}

\section{تأويل (وإذا تولى سعى في الأرض ليفسد فيها)}
فسر \rawi{الطبري} هذا السعي بأنه العمل في الأرض بالمعصية والإفساد بعد الإدبار عن مجلس النبي صلى الله عليه وسلم أو بعد التمكن من الفساد. ثم ذكر أن الإفساد يشمل صوراً متعددة:
\begin{itemize}
    \item كقطع الطريق.
    \item وترويع السبل.
    \item وقطع الأرحام.
    \item وسفك الدماء.
    \item وإهلاك الحرث والنسل.
\end{itemize}

ورجح أن اللفظ عام في كل فساد، مع ملاحظة أن ظاهر السياق يقوي بعض صوره أكثر من بعض.

\section{تأويل (ومن الناس من يشري نفسه ابتغاء مرضات الله)}
جاءت هذه الآية في مقابلة الصنف الأول، فوصفت من يبيع نفسه لله ابتغاء مرضاته، سواء كان ذلك في الجهاد، أو الهجرة، أو الأمر بالمعروف والنهي عن المنكر، أو بذل النفس في طاعة الله على أي وجه دخل في ظاهر الآية.

ونقل \rawi{الطبري} ما ورد في نزولها في صهيب وغيره، ثم ما ورد عن \rawi{عمر} و\rawi{ابن عباس} وغيرهما في حملها على من يأمر بالمعروف وينهى عن المنكر ويحتسب نفسه في مقاومة أهل الفساد.

ورجح أن ظاهر الآية يعم كل من باع نفسه لله في طاعته، وأن خصوص السبب لا يمنع عموم الوصف.

\begin{fawaid}[أعظم مقاومة للمنكر نشر المعروف وتيسير سبله]
ليس النهي عن المنكر مقصوراً على مواجهة فاعله مباشرة، بل من أعظم صوره نشر المعروف، وفتح أبوابه، وشغل الناس بالحق والنافع؛ لأن هذا من أبلغ ما يضعف الباطل ويصرف القلوب عنه.
\end{fawaid}

\begin{fawaid}[من رأفة الله بعباده أن يقيم فيهم من يبيع نفسه لنصرة الحق]
ختمت الآية بقوله تعالى \begin{ayah}وَاللَّهُ رَءُوفٌ بِالْعِبَادِ\end{ayah}، ومن معاني ذلك أن من رأفته بعباده أن يوجد فيهم هذا الصنف الشريف الذي يحتسب نفسه لله فيجاهد، ويهاجر، ويأمر بالمعروف، ويدفع أهل الفساد.
\end{fawaid}

\separator

ختم الشيخ اللقاء بالتأكيد على غنى القرآن والسنة عن كثير مما يطلبه الناس في المراجع المتأخرة في أبواب التربية والإصلاح والوعظ؛ فكلما ازداد العبد صلة بكتاب الله وهدي رسوله صلى الله عليه وسلم استغنى بهما، ورأى فيهما المنطلق والميزان والمرد في أبواب الدين كلها.
